\documentclass[11pt,a4paper]{article}
\usepackage[T1]{fontenc}
\usepackage[utf8]{inputenc}
\usepackage{lmodern}
\usepackage{amsmath,amssymb,amsthm,mathtools,mathrsfs}
\usepackage{graphicx,float,xcolor,multicol}
\usepackage[shortlabels]{enumitem}
\usepackage[margin=27mm]{geometry}
\usepackage{microtype,needspace,placeins}
\usepackage{tikz,tikz-cd}
\usetikzlibrary{arrows.meta,calc,positioning,patterns,decorations.pathreplacing}
\usepackage[colorlinks=true,linkcolor=teal,urlcolor=teal,citecolor=teal]{hyperref}
\setlength{\parindent}{0pt}
\setlength{\parskip}{.6em}
\setcounter{secnumdepth}{0}
\allowdisplaybreaks
\raggedbottom
\providecommand{\cross}{\times}
\DeclareUnicodeCharacter{2020}{\ensuremath{\dagger}}
\DeclareMathOperator{\supp}{supp}
\DeclareMathOperator*{\esssup}{ess\,sup}
\providecommand{\chapter}{\section}
\newenvironment{tool}[2]{\subsection*{Tool #1 --- #2}}{}
\newcommand{\ToolRef}[1]{Tool #1}
\newcommand{\FigureTag}[2]{\textit{Figure #1}}
\newcommand{\Result}[1]{\subsection*{#1}}
\newcommand{\Application}[1]{\subsubsection*{Application\if\relax\detokenize{#1}\relax\else\space--- #1\fi}}
\newcommand{\ProofHeading}{\par\textit{Proof.}\space}
\newcommand{\ProofOf}[1]{\par\textit{Proof of #1.}\space}
\newcommand{\RemarkHeading}{\par\textbf{Remark.}\space}
\newcommand{\NoteHeading}{\par\textbf{Note.}\space}
\newcommand{\ExerciseHeading}{\par\textbf{Exercise.}\space}

% Rendering support only; mathematical text is in each independent post.
\usepackage{booktabs,array,longtable}
\definecolor{HandoutBlue}{HTML}{2F6690}
\definecolor{HandoutNavy}{HTML}{17365D}
\newtheorem{theorem}{Theorem}
\newtheorem{lemma}[theorem]{Lemma}
\newtheorem{proposition}[theorem]{Proposition}
\newtheorem{corollary}[theorem]{Corollary}
\newtheorem{conjecture}[theorem]{Conjecture}
\newtheorem{definition}{Definition}
\newtheorem{example}{Example}
\newtheorem{namedtoolcounter}{Tool}
\newenvironment{namedtool}[1]{\begin{namedtoolcounter}[#1]}{\end{namedtoolcounter}}
\newtheorem*{claim}{Claim}
\newtheorem*{remark}{Remark}
\newtheorem*{exercise}{Exercise}
\newtheorem*{question}{Question}
\newtheorem*{observation}{Observation}
\newcommand{\SourceChapter}[2]{\section*{#2}}
\newcommand{\Lecture}[2]{\section*{Lecture #1: #2}}
\newcommand{\unclear}[1]{\textit{[unclear: #1]}}
\newcommand{\sourcegap}{\textit{[unfinished in the source]}}
\DeclareMathOperator{\dist}{dist}
\DeclareMathOperator{\id}{id}
\DeclareMathOperator{\Hol}{Hol}
\DeclareMathOperator{\Per}{Per}
\DeclareMathOperator{\Fix}{Fix}
\DeclareMathOperator{\diam}{diam}
\DeclareMathOperator{\Var}{Var}
\DeclareMathOperator{\Leb}{Leb}
\DeclareMathOperator{\Hom}{Hom}
\DeclareMathOperator{\End}{End}
\DeclareMathOperator{\Ann}{Ann}
\DeclareMathOperator{\Spec}{Spec}
\DeclareMathOperator{\Max}{Max}
\DeclareMathOperator{\Ob}{Ob}
\DeclareMathOperator{\Tor}{Tor}
\DeclareMathOperator{\Ass}{Ass}
\DeclareMathOperator{\Mor}{Mor}
\DeclareMathOperator{\Fun}{Fun}
\DeclareMathOperator{\Nat}{Nat}
\DeclareMathOperator{\Cone}{Cone}
\DeclareMathOperator{\MCG}{MCG}
\DeclareMathOperator{\Aut}{Aut}
\DeclareMathOperator{\Deck}{Deck}
\DeclareMathOperator{\SL}{SL}
\DeclareMathOperator{\PSL}{PSL}
\DeclareMathOperator{\Area}{Area}
\DeclareMathOperator{\im}{im}
\DeclareMathOperator{\PSh}{PSh}
\newcommand{\CRing}{\mathsf{CRings}}
\newcommand{\AMod}{A\text{-}\mathsf{Mod}}
\newcommand{\Ab}{\mathsf{Ab}}
\newcommand{\Z}{\mathbb Z}
\newcommand{\Q}{\mathbb Q}
\newcommand{\R}{\mathbb R}
\newcommand{\C}{\mathbb C}
\newcommand{\N}{\mathbb N}
\newcommand{\kk}{\mathbb k}
\renewcommand{\aa}{\mathfrak a}
\newcommand{\bb}{\mathfrak b}
\newcommand{\pp}{\mathfrak p}
\newcommand{\qq}{\mathfrak q}
\newcommand{\mm}{\mathfrak m}
\newcommand{\nn}{\mathfrak n}
\newcommand{\rad}{\mathrm r}
\newcommand{\cat}[1]{\mathsf{#1}}
\setcounter{secnumdepth}{0}
\allowdisplaybreaks

\graphicspath{{figures/lemma-book/}}
\begin{document}
\section*{Thrilling Thirty: The Surviving Problems}
% BLOG-CONTENT-BEGIN
\subsection{A note to the reader}
This is not the end of our journey; there is a lot more to come.
I decided not to write the 100th lemma. Rather, I have given you, the reader,
the option to choose one of the best ideas and title that the 100th lemma!
Anyway, I myself formed a thrilling thirty problems. Take that as a sweet dessert
after such a memorable feast.

\emph{Happy Problem Solving!}

Do not forget: I'll be back. Ba Bum Bum Bum\ldots


\subsection{Thrilling Thirty}
\begin{enumerate}[label=\textbf{\arabic*.},itemsep=.75em]
\setcounter{enumi}{0}\item In $\triangle ABC$, $M$ is the midpoint of the base $BC$. The triangle is
dissected into two parts of equal area. Find a way to dissect one piece into as few
pieces as possible so that their rearrangement tiles the other exactly.


\setcounter{enumi}{6}\item Find all pairs $(n,x)$, where $n\in\mathbb N$ and $0\leq x\leq2$, such that
$\sum_{k=1}^{n}\arcsin(x/2^k)=\pi$.

% Source page 65
\setcounter{enumi}{7}\item Prove that for every $n$ there exists an $m$ such that $m^4\equiv406\pmod{3^n}$.


\setcounter{enumi}{9}\item Suppose $1/2014=1/x_1+1/x_2+\cdots+1/x_n$, where
$x_n>x_{n-1}>\cdots>x_2>x_1>1$ are positive integers. Find $n$ and the $x_i$ satisfying
this equation. Are there infinitely many solutions?


\setcounter{enumi}{11}\item Compute the following ascending fraction:
\[
1+\cfrac{1+\cfrac{1+\cfrac{1+\cfrac{1+\cfrac{1+\cdots}{4870847}}{2207}}{47}}{7}}{3}.
\]
Here $7=3^2-2$, $47=7^2-2$, $2207=47^2-2$, $4870847=2207^2-2$, and so on.



\setcounter{enumi}{18}

\setcounter{enumi}{19}\item Triangle $ABC$ is equilateral, $P$ is a point on its circumcircle, and $Q$ is
any point outside the circle such that $PB=PQ=5\,\mathrm m$. Vedant and Rahul start
running along $PA$ and $PB$, respectively, with the same speed $v$. I start running
along $QP$ with speed $u$; after reaching $P$, I decide to run with speed $v$ towards $C$.
\par It is found that Vedant and I take the same time to reach our respective destinations.
Vedant says that we all run uniformly on our paths. Is Vedant a genius to point that
out, or was he just out of his mind!
\par Furthermore, Rahul and I inform Vedant that it took us 5 and 10 seconds,
respectively, to reach our destinations. Vedant responds that if $BC=BQ$, then
$\triangle BQC$ is equilateral. Is Vedant a brainbox, or are we dumbasses?

\setcounter{enumi}{20}\item Imagine a frictionless rectangular billiard table with dimensions
$(503\times p)$ metres, where $p$ is prime. A ball is hit from a corner of the table
with velocity $\sqrt{2p^2}\,\mathrm{m\,s^{-1}}$, making an angle of $45^\circ$ with the rails.
The ball rolls on the table for about 503 seconds before falling into a pocket.
Vedant miscounts the number of times the ball hits a wall of the table, getting four
more than the actual number. What number does Vedant obtain if the ball lands in
the pocket from which it was originally shot?

\setcounter{enumi}{21}\item Vedant discovers a new particle. To respect his contribution, let us name it
``VEDU'' in his honour! It has a weird property: two particles brought close together
bind to form a single particle.
\par There are three mutually intersecting unit circles such that the system of
circles is symmetric about the lines $PP^*,QQ^*,RR^*$, where these are their respective
common chords. Let $A,B,C$ be the midpoints of arcs $PR,RQ,PQ$, respectively, such that
the area of $\triangle ABC$ is maximum. Six VEDUs are launched simultaneously with
the same velocity along the circular paths $AP,AR,BQ,BR,CQ,CP$. If the length of arc
$PR$ is $x$, determine the velocity of each particle so that exactly after 2014 seconds
all that remains is one VEDU.

% Source page 68 continues Problem 22 above.
\setcounter{enumi}{22}\item Prove that one cannot dissect a regular pentagon into exactly 2014 triangular
pieces by a series of consecutive cuts. Furthermore, prove that one can always
divide it into 2013 triangular pieces.
\par A cut is a straight line segment connecting two vertices. A new vertex is formed
where two or more cuts intersect. A cut may connect two new vertices, two original
vertices, or an original vertex and a new vertex.

\setcounter{enumi}{23}\item Obama and Osama are separated by some distance. A mirror starts moving from
Obama to Osama along a certain path such that they can see each other. They spend
time teasing each other (of course, by making those bullshit faces). Throughout the
mirror's journey, the sum of the angles of incidence and reflection of the light
rays emanating from Obama equals the total angle of rotation of the mirror about
the point of incidence $P$. During the journey, there is a stage at which $P$ is
equidistant from Obama and Osama. If that distance is $20\sqrt7$ metres, find the
distance between Obama and Osama.

\setcounter{enumi}{24}\item Given positive reals $a,b,c$, prove
\[
\sum_{\mathrm{cyc}}\frac{(a+1)^2+(b+1)^2}{(a+b)(a+b+c)}>\frac{2014}{1024}.
\]

\setcounter{enumi}{25}\item Given nine distinct integers satisfying
$ (x_1^2+x_2^2+x_3^2)+(y_1^2+y_2^2+y_3^2)=z_1^2+z_2^2+z_3^2$,
prove that there are infinitely many families of solutions of this Diophantine
equation, even if no three of the nine integers
$(x_1,x_2,x_3,y_1,y_2,y_3,z_1,z_2,z_3)$ form a Pythagorean triplet.
\par\emph{Open problem:} Can one find all families of solutions of this Diophantine equation?

% Source page 69
\setcounter{enumi}{26}\item Elementary puzzle:
\begin{itemize}
\item If you take away my last two digits and place them in any order in front of me,
I am an $x$-digit number, where $x$ is my first digit.
\item If you take the first digit raised to the reciprocal of the number of sevens
in me, you get my seventh digit.
\item Take the difference of my seventh and eighth digits: you get the number of
zeroes in me. If you take the sum, you get back my first digit.
\item If you remove my first digit, my number of digits equals my seventh digit
from the last. I behave the same way if you remove my first two digits.
\end{itemize}

\setcounter{enumi}{27}\item For every nonzero integer $n$, there exists a pair of primes $(p,q)$ such that
$n\mid(p-q)$. Such pairs also occur infinitely often.

\end{enumerate}
\begin{center}\sffamily\bfseries\color{HandoutBlue}THE END??\end{center}
% BLOG-CONTENT-END
\end{document}
